\chapter{Dérivation~: un premier cours}\label{ch:g11:deriv}

À quelle vitesse une \omterm{def:g11:func:function}{fonction} change-t-elle~? La réponse en un point
unique est la \emph{\omterm{def:g11:deriv:derivative}{dérivée}}, le nombre le plus utile de toutes les
mathématiques appliquées. Ce chapitre la construit à partir des \omterm{def:g11:deriv:rate}{taux de
variation moyens}, la calcule pour les \omterm{def:g11:func:reference}{fonctions de référence}, et utilise
son signe pour lire les variations d'une \omterm{def:g11:func:function}{fonction}. Tout ce qui suit est
étendu et approfondi dans le \cref{ch:g12:deriv}.

\section{Le taux de variation moyen}

\begin{definition}[Taux de variation moyen]\label{def:g11:deriv:rate}
Soit $f$ définie sur un \omterm{def:g10:numbers:interval}{intervalle} $I$, et $a \neq b$ dans $I$. Le
\emph{taux de variation moyen}\index{taux de variation moyen} de $f$
entre $a$ et $b$ est
\[
\frac{f(b) - f(a)}{b - a}.
\]
C'est le \omterm{def:g10:reffunc:affine}{coefficient directeur} de la \emph{sécante}\index{sécante} passant par les
points $\bigl(a, f(a)\bigr)$ et $\bigl(b, f(b)\bigr)$ de la courbe.
\end{definition}

\begin{example}
Une voiture parcourt $140$~km entre 14~h et 16~h~: sa vitesse \omterm{def:g10:stats:mean}{moyenne} est
$\frac{140}{2} = 70$~km/h. Le compteur, en revanche, affiche une vitesse
\emph{instantanée} à chaque instant --- c'est exactement la notion que
la \omterm{def:g11:deriv:derivative}{dérivée} rend précise.
\end{example}

\section{La dérivée en un point}

\begin{definition}[Dérivée en un point]\label{def:g11:deriv:derivative}
Soit $f$ définie sur un \omterm{def:g10:numbers:interval}{intervalle} contenant $a$. Pour $h \neq 0$, former
le \omterm{def:g11:deriv:rate}{taux de variation moyen} entre $a$ et $a + h$~:
\[
r(h) = \frac{f(a+h) - f(a)}{h}.
\]
Si $r(h)$ se rapproche d'un unique nombre fixe lorsque $h$ se rapproche
arbitrairement de $0$, on dit que $f$ est \emph{dérivable en
$a$}\index{dérivée}\index{dérivable}~; ce nombre est la \emph{dérivée de $f$ en $a$},
notée $f'(a)$~:
\[
f'(a) = \lim_{h \to 0} \frac{f(a+h) - f(a)}{h}.
\]
\end{definition}

\begin{example}\label{ex:g11:deriv:atpoint}
Soit $f(x) = x^2$ et $a = 1$. Pour $h \neq 0$~:
\[
\frac{f(1+h) - f(1)}{h}
= \frac{(1+h)^2 - 1}{h}
= \frac{1 + 2h + h^2 - 1}{h}
= \frac{h(2 + h)}{h}
= 2 + h .
\]
Lorsque $h$ tend vers $0$, cela tend vers $2$~: la \omterm{def:g11:func:function}{fonction} $x^2$ est
\omterm{def:g11:deriv:derivative}{dérivable} en $1$ et $f'(1) = 2$. Noter la stratégie~: \emph{simplifier
le quotient jusqu'à ce que $h$ n'apparaisse plus au dénominateur}, puis
faire $h \to 0$.
\end{example}

\begin{example}[Un point non dérivable]
Soit $f(x) = \abs{x}$ et $a = 0$. Alors
$\frac{f(h) - f(0)}{h} = \frac{\abs h}{h}$, qui vaut $1$ pour $h > 0$ et
$-1$ pour $h < 0$~: aucun nombre unique n'est approché, et $\abs x$
n'est pas \omterm{def:g11:deriv:derivative}{dérivable} en $0$. Sa courbe a un \emph{coin} là.
\end{example}

\section{La tangente}

\begin{definition}[Tangente]\label{def:g11:deriv:tangent}
Si $f$ est \omterm{def:g11:deriv:derivative}{dérivable} en $a$, la \emph{tangente}\index{tangente} à la
courbe au point $A = \bigl(a, f(a)\bigr)$ est la droite passant par $A$
de \omterm{def:g10:reffunc:affine}{coefficient directeur} $f'(a)$. Son \omterm{def:g10:algebra:equation}{équation} est
\[
y = f(a) + f'(a)(x - a).
\]
\end{definition}

La \omterm{def:g11:deriv:tangent}{tangente} est la position limite des \omterm{def:g11:deriv:rate}{sécantes} par $A$~: lorsque
$h \to 0$, le second point $\bigl(a+h, f(a+h)\bigr)$ glisse le long de la
courbe vers $A$, et le \omterm{def:g10:reffunc:affine}{coefficient directeur} de la \omterm{def:g11:deriv:rate}{sécante}
$\frac{f(a+h)-f(a)}{h}$ tend vers $f'(a)$.

\begin{omfigure}
\begin{tikzpicture}
\begin{axis}[omaxis,
  width=8.8cm, height=6cm,
  xmin=-0.4, xmax=3.2, ymin=-0.6, ymax=4.3,
  xtick={1,2.2}, xticklabels={$a$,$a+h$}, ytick=\empty]
  \addplot[omProp, thick, domain=-0.1:2.9] {1 + 1.6*(x-1)};
  \addplot[omThm, thick, domain=-0.35:2.6] {1 + 1*(x-1)};
  \addplot[omDef, thick, domain=-0.4:2.75] {x^2/2 + 0.5};
  \fill (axis cs:1,1) circle (1.5pt) node[above left, font=\small] {$A$};
  \fill (axis cs:2.2,2.92) circle (1.5pt)
    node[above left, font=\small] {$B$};
  \node[omProp, font=\small, anchor=west] at (axis cs:2.62,3.75)
    {sécante};
  \node[omThm, font=\small, anchor=west] at (axis cs:2.62,2.45)
    {tangente};
\end{axis}
\end{tikzpicture}

{\small La \omterm{def:g11:deriv:rate}{sécante} par $A$ et $B = \bigl(a+h, f(a+h)\bigr)$ a pour
\omterm{def:g10:reffunc:affine}{coefficient directeur} $\frac{f(a+h)-f(a)}{h}$. Lorsque $B$ glisse vers
$A$, la \omterm{def:g11:deriv:rate}{sécante} bascule dans la \omterm{def:g11:deriv:tangent}{tangente}, dont le \omterm{def:g10:reffunc:affine}{coefficient directeur}
est $f'(a)$.}
\end{omfigure}

\begin{example}\label{ex:g11:deriv:tangenteq}
Pour $f(x) = x^2$ en $a = 1$~: $f(1) = 1$ et $f'(1) = 2$
(\cref{ex:g11:deriv:atpoint}), donc la \omterm{def:g11:deriv:tangent}{tangente} est
\[
y = 1 + 2(x - 1) = 2x - 1 .
\]
C'est la droite $d_2$ trouvée par pure algèbre dans
l'\cref{exo:g11:quad:11}.
\end{example}

\section{Dérivées des fonctions de référence}

\begin{proposition}\label{prop:g11:deriv:refderivatives}
Sur leurs ensembles de dérivabilité~:
\[
(c)' = 0, \quad
(x)' = 1, \quad
(x^2)' = 2x, \quad
(x^3)' = 3x^2, \quad
\left(\frac1x\right)' = -\frac{1}{x^2}, \quad
(\sqrt x\,)' = \frac{1}{2\sqrt x} \ (x > 0).
\]
Plus généralement, $(x^n)' = n\,x^{n-1}$ pour tout entier $n \geq 1$.
\end{proposition}

\begin{proof}
Chaque démonstration suit la stratégie de
l'\cref{ex:g11:deriv:atpoint}~: simplifier le taux, puis faire
$h \to 0$. Fixer $a$ dans l'\omterm{def:g11:func:function}{ensemble de définition}.

\emph{Constante et identité.} $\frac{c - c}{h} = 0$ et
$\frac{(a+h) - a}{h} = 1$ pour tout $h$~: les limites sont $0$ et $1$.

\emph{Carré.} $\frac{(a+h)^2 - a^2}{h} = \frac{2ah + h^2}{h} = 2a + h
\to 2a$.

\emph{Cube.} $(a+h)^3 = a^3 + 3a^2h + 3ah^2 + h^3$, donc le taux est
$3a^2 + 3ah + h^2 \to 3a^2$.

\emph{Inverse.} Pour $a \neq 0$ et $h$ assez petit pour que
$a + h \neq 0$~:
\[
\frac{1}{h}\left(\frac{1}{a+h} - \frac{1}{a}\right)
= \frac{1}{h} \cdot \frac{a - (a+h)}{a(a+h)}
= \frac{-1}{a(a+h)}
\xrightarrow[h \to 0]{} -\frac{1}{a^2}.
\]

\emph{Racine carrée.} Pour $a > 0$, multiplier par le conjugué~:
\[
\frac{\sqrt{a+h} - \sqrt a}{h}
= \frac{(a + h) - a}{h\,(\sqrt{a+h} + \sqrt a)}
= \frac{1}{\sqrt{a+h} + \sqrt a}
\xrightarrow[h \to 0]{} \frac{1}{2\sqrt a}.
\]

\emph{Puissance générale.} Le motif $2a$, $3a^2$ se poursuit~: en
développant $(a+h)^n$, les termes après $a^n + n\,a^{n-1}h$ portent tous
$h^2$, donc le taux tend vers $n\,a^{n-1}$. (Une induction complète est
donnée dans le \cref{ch:g12:deriv}.)
\end{proof}

\section{Règles de dérivation}

\begin{theorem}[Opérations]\label{thm:g11:deriv:rules}
Soient $u$ et $v$ \omterm{def:g11:deriv:derivative}{dérivables} sur $I$ et $\lambda \in \R$. Alors
$u + v$, $\lambda u$, $uv$ sont \omterm{def:g11:deriv:derivative}{dérivables} sur $I$, de même que
$\frac{u}{v}$ là où $v \neq 0$, et
\[
(u+v)' = u' + v', \qquad
(\lambda u)' = \lambda u', \qquad
(uv)' = u'v + uv', \qquad
\left(\frac{u}{v}\right)' = \frac{u'v - uv'}{v^2}.
\]
\end{theorem}

\begin{proof}[Démonstration des trois premières règles]
Fixer $a \in I$ et abréger $\Delta u = u(a+h) - u(a)$,
$\Delta v = v(a+h) - v(a)$.

\emph{Somme.} Le taux de $u + v$ est
$\frac{\Delta u + \Delta v}{h} = \frac{\Delta u}{h} + \frac{\Delta v}{h}
\to u'(a) + v'(a)$.

\emph{Multiple.} Le taux de $\lambda u$ est
$\lambda \frac{\Delta u}{h} \to \lambda u'(a)$.

\emph{Produit.} Ajouter et soustraire le terme mixte $u(a+h)\,v(a)$~:
\begin{align*}
\frac{u(a+h)v(a+h) - u(a)v(a)}{h}
&= \frac{u(a+h)v(a+h) - u(a+h)v(a) + u(a+h)v(a) - u(a)v(a)}{h}\\
&= u(a+h)\,\frac{\Delta v}{h} + v(a)\,\frac{\Delta u}{h}.
\end{align*}
Lorsque $h \to 0$~: $\frac{\Delta v}{h} \to v'(a)$,
$\frac{\Delta u}{h} \to u'(a)$, et $u(a+h) \to u(a)$ (son taux a une
limite, donc $\Delta u = h \cdot \frac{\Delta u}{h} \to 0$). Le taux
tend donc vers $u(a)v'(a) + v(a)u'(a)$.

La règle du quotient se démontre de la même façon~; elle est
\emph{admise à ce niveau} et prouvée dans le \cref{ch:g12:deriv}.
\end{proof}

\begin{example}\label{ex:g11:deriv:computations}
\emph{Polynôme.} $f(x) = 3x^3 - 5x^2 + 7x - 2$~:
\[
f'(x) = 3 \times 3x^2 - 5 \times 2x + 7 = 9x^2 - 10x + 7 .
\]

\emph{Quotient.} $g(x) = \frac{2x+1}{x^2+1}$~: avec $u = 2x + 1$,
$v = x^2 + 1$,
\[
g'(x) = \frac{2(x^2+1) - (2x+1)\,2x}{(x^2+1)^2}
= \frac{-2x^2 - 2x + 2}{(x^2+1)^2}.
\]

\emph{Produit.} $k(x) = x\sqrt x$ pour $x > 0$~:
$k'(x) = 1 \times \sqrt x + x \times \frac{1}{2\sqrt x}
= \sqrt x + \frac{\sqrt x}{2} = \frac{3\sqrt x}{2}$.
\end{example}

\section{Signe de la dérivée et variations}

\begin{theorem}[Dérivée et variations]\label{thm:g11:deriv:variations}
Soit $f$ \omterm{def:g11:deriv:derivative}{dérivable} sur un \omterm{def:g10:numbers:interval}{intervalle} $I$.
\begin{enumerate}
  \item Si $f' \geq 0$ sur $I$, alors $f$ est \omterm{def:g11:func:monotone}{croissante} sur $I$.
  \item Si $f' \leq 0$ sur $I$, alors $f$ est \omterm{def:g10:functions:variations}{décroissante} sur $I$.
  \item Si $f' > 0$ sur $I$ (sauf éventuellement en un nombre fini de
        points où elle s'annule), alors $f$ est strictement \omterm{def:g11:func:monotone}{croissante}
        sur $I$.
\end{enumerate}
\end{theorem}

\begin{proof}
L'intuition est claire --- une courbe dont toutes les \omterm{def:g11:deriv:tangent}{tangentes} pointent
vers le haut doit monter --- mais une démonstration rigoureuse nécessite
le théorème des accroissements finis. Le résultat est \emph{admis à ce
niveau}~; voir le \cref{thm:g12:deriv:variations}.
\end{proof}

\begin{proposition}[Extrema]\label{prop:g11:deriv:extrema}
Si $f'$ change de signe en $a$ de $+$ à $-$, alors $f$ a un \omterm{def:g10:functions:extrema}{maximum}
local en $a$~; de $-$ à $+$, un \omterm{def:g10:functions:extrema}{minimum} local. En un tel point,
$f'(a) = 0$ et la \omterm{def:g11:deriv:tangent}{tangente} est horizontale.
\end{proposition}

\begin{proof}
Si $f' \geq 0$ sur un \omterm{def:g10:numbers:interval}{intervalle} se terminant en $a$ et $f' \leq 0$ sur
un \omterm{def:g10:numbers:interval}{intervalle} commençant en $a$, alors $f$ croît jusqu'à $f(a)$ et
décroît ensuite (\cref{thm:g11:deriv:variations})~: $f(a)$ est le plus
grand au voisinage. Comme $f'$ change de signe en $a$ et y est définie,
$f'(a) = 0$.
\end{proof}

\begin{method}[Étudier une fonction]\label{met:g11:deriv:study}
\begin{enumerate}
  \item Calculer $f'(x)$ et le simplifier --- idéalement sous forme
        factorisée.
  \item Déterminer le signe de $f'(x)$ sur chaque \omterm{def:g10:numbers:interval}{intervalle} (pour un
        $f'$ du second degré, utiliser le \cref{thm:g11:quad:sign}).
  \item Enregistrer les résultats dans un \emph{\omterm{def:g10:functions:table}{tableau de variations}}~:
        une ligne pour le signe de $f'$, une ligne de flèches pour $f$,
        avec les valeurs de $f$ aux points de retournement.
  \item Conclure~: extrema, nombre de solutions de $f(x) = k$,
        inégalités.
\end{enumerate}
\end{method}

\begin{example}\label{ex:g11:deriv:study}
Étudier $f(x) = x^3 - 3x + 1$ sur $\R$. D'abord,
$f'(x) = 3x^2 - 3 = 3(x-1)(x+1)$~: un trinôme de \omterm{def:g11:quad:discriminant}{racines} $-1$ et $1$,
positif à l'extérieur. Donc $f$ est \omterm{def:g11:func:monotone}{croissante} sur
$\intoc{-\infty}{-1}$, \omterm{def:g10:functions:variations}{décroissante} sur $\intcc{-1}{1}$, \omterm{def:g11:func:monotone}{croissante} sur
$\intco{1}{+\infty}$, avec un \omterm{def:g10:functions:extrema}{maximum} local $f(-1) = -1 + 3 + 1 = 3$ et
un \omterm{def:g10:functions:extrema}{minimum} local $f(1) = 1 - 3 + 1 = -1$.
\end{example}

\begin{omfigure}
\begin{tikzpicture}
\begin{axis}[omaxis,
  width=8.4cm, height=5.8cm,
  xmin=-2.6, xmax=2.6, ymin=-3.6, ymax=4.4,
  xtick={-1,1}, ytick={-1,3}]
  \addplot[omProp, thick, domain=-1.7:-0.3] {3};
  \addplot[omProp, thick, domain=0.3:1.7] {-1};
  \addplot[omDef, thick, domain=-2.25:2.25, samples=90] {x^3 - 3*x + 1};
  \fill (axis cs:-1,3) circle (1.5pt);
  \fill (axis cs:1,-1) circle (1.5pt);
\end{axis}
\end{tikzpicture}

{\small La courbe de $f(x) = x^3 - 3x + 1$
(\cref{ex:g11:deriv:study})~: \omterm{def:g11:deriv:tangent}{tangentes} horizontales (orange) au
\omterm{def:g10:functions:extrema}{maximum} local $(-1, 3)$ et au \omterm{def:g10:functions:extrema}{minimum} local $(1, -1)$, où $f'$ change de
signe.}
\end{omfigure}

\begin{example}[Optimisation]\label{ex:g11:deriv:box}
Une boîte ouverte est fabriquée à partir d'une feuille de
$12 \times 12$~cm en découpant un carré de côté $x$ à chaque coin et en
pliant. Son volume est $V(x) = x(12 - 2x)^2$ pour $0 < x < 6$. En
développant, $V(x) = 4x^3 - 48x^2 + 144x$, donc
\[
V'(x) = 12x^2 - 96x + 144 = 12(x^2 - 8x + 12) = 12(x - 2)(x - 6).
\]
Sur $\intoo{0}{6}$, le facteur $x - 6$ est négatif, donc $V'(x)$ a le
signe de $-(x-2)$~: positif avant $2$, négatif après. Le volume est
maximal pour $x = 2$~cm, donnant $V(2) = 2 \times 8^2 = 128$~cm$^3$.
\end{example}

\section{Exercices}

\begin{exercise}[$\star$]\label{exo:g11:deriv:1}
En utilisant la définition (taux de variation, puis $h \to 0$), calculer
$f'(2)$ pour $f(x) = x^2$, puis $g'(1)$ pour $g(x) = x^2 + 3x$.
\end{exercise}

\begin{exercise}[$\star$]\label{exo:g11:deriv:2}
Dériver~:
\[
f(x) = 4x^3 - 2x^2 + x - 7, \qquad
g(x) = 5\sqrt{x} + \frac{3}{x}, \qquad
h(x) = (2x+1)(x^2 - 3) .
\]
\end{exercise}

\begin{exercise}[$\star$]\label{exo:g11:deriv:3}
Donner l'\omterm{def:g10:algebra:equation}{équation} de la \omterm{def:g11:deriv:tangent}{tangente} à la courbe de $f$ au point d'\omterm{def:g10:coordgeom:system}{abscisse}
$a$~:
\[
f(x) = x^2 - 3x, \ a = 2; \qquad
f(x) = \frac1x, \ a = 1; \qquad
f(x) = \sqrt x, \ a = 4 .
\]
\end{exercise}

\begin{exercise}[$\star$]\label{exo:g11:deriv:4}
Dériver les quotients~:
\[
f(x) = \frac{x+2}{x-1}, \qquad
g(x) = \frac{x^2}{x^2+1}, \qquad
h(x) = \frac{1}{x^2 + x + 1} .
\]
\end{exercise}

\begin{exercise}[$\star\star$]\label{exo:g11:deriv:5}
Étudier les variations de $f(x) = x^3 - 6x^2 + 9x - 2$ sur $\R$ (\omterm{def:g10:functions:table}{tableau
de variations}, extrema locaux).
\end{exercise}

\begin{exercise}[$\star\star$]\label{exo:g11:deriv:6}
Étudier les variations de $f(x) = \frac{x^2 + 3}{x + 1}$ sur
$\intoo{-1}{+\infty}$.
\end{exercise}

\begin{exercise}[$\star\star$]\label{exo:g11:deriv:7}
Soit $f(x) = x^2 - 4x + 5$.
\begin{enumerate}
  \item Trouver le point de la courbe où la \omterm{def:g11:deriv:tangent}{tangente} est horizontale.
  \item Trouver le ou les points où la \omterm{def:g11:deriv:tangent}{tangente} est parallèle à la
        droite $y = 2x + 1$.
\end{enumerate}
\end{exercise}

\begin{exercise}[$\star\star$]\label{exo:g11:deriv:8}
Un enclos rectangulaire est construit contre un mur avec $60$~m de
clôture (trois côtés). Exprimer l'aire en \omterm{def:g11:func:function}{fonction} de la largeur $x$, et
utiliser la \omterm{def:g11:deriv:derivative}{dérivée} pour trouver les dimensions d'aire maximale.
(Comparer avec la méthode du sommet du \cref{ch:g11:quad}.)
\end{exercise}

\begin{exercise}[$\star\star$]\label{exo:g11:deriv:9}
Montrer que pour tout $x > 0$, $\sqrt{x} \leq \frac{x + 1}{2}$, en
étudiant la \omterm{def:g11:func:function}{fonction} $f(x) = \frac{x+1}{2} - \sqrt x$ sur
$\intoo{0}{+\infty}$.
\end{exercise}

\begin{exercise}[$\star\star$]\label{exo:g11:deriv:10}
Une boîte cylindrique doit contenir $500$~cm$^3$. Son aire totale (deux
disques plus le côté) est $S = 2\pi r^2 + 2\pi r h$ avec
$\pi r^2 h = 500$. Exprimer $S$ en \omterm{def:g11:func:function}{fonction} de $r$ seul, et montrer que
$S'(r) = 4\pi r - \frac{1000}{r^2}$. En déduire le rayon qui minimise
l'aire, sous forme exacte.
\end{exercise}

\begin{exercise}[$\star\star\star$]\label{exo:g11:deriv:11}
Combien de \omterm{def:g11:deriv:tangent}{tangentes} à la \omterm{def:g11:quad:parabola}{parabole} $y = x^2$ passent par le point
extérieur $P(1, -3)$~? Déterminer leurs \omterm{def:g10:algebra:equation}{équations}. (Soit $a$ l'\omterm{def:g10:coordgeom:system}{abscisse}
du point de contact et exprimer que la \omterm{def:g11:deriv:tangent}{tangente} en $a$ passe par $P$.)
\end{exercise}

\begin{exercise}[$\star\star\star$]\label{exo:g11:deriv:12}
Soit $f(x) = x^3 - 3x + 1$ (voir l'\cref{ex:g11:deriv:study}). À l'aide
du \omterm{def:g10:functions:table}{tableau de variations}, déterminer le nombre de solutions de
l'\omterm{def:g10:algebra:equation}{équation} $f(x) = k$ selon la valeur du réel $k$.
\end{exercise}

\section{Problème~: les trois métiers de la tangente}

\begin{problem}[{Devoir du week-end --- le miroir parabolique
démontré, la machine à trouver les racines de Newton, et la poutre la
plus résistante du tronc}]
\label{pb:g11:deriv:1}
Une \omterm{def:g11:deriv:derivative}{dérivée} trace des \omterm{def:g11:deriv:tangent}{tangentes}~; le talent paraît modeste. Ce
problème montre la \omterm{def:g11:deriv:tangent}{tangente} exerçant trois métiers à la fois~: elle
démontre le théorème du phare laissé en suspens au
\cref{pb:g11:quad:1} (tout rayon parallèle à l'axe d'une \omterm{def:g11:quad:parabola}{parabole} se
réfléchit en passant par le \omterm{pb:g11:quad:1}{foyer}), elle fait tourner l'algorithme de
résolution d'\omterm{def:g10:algebra:equation}{équations} le plus rapide en usage courant --- présent
dans chaque calculatrice et dans chaque réseau de neurones --- et
elle trouve la poutre rectangulaire la plus résistante qu'un tronc
rond puisse fournir.

\medskip
\noindent\textbf{Partie I --- Aisance avec les \omterm{def:g11:deriv:tangent}{tangentes}.}
\begin{enumerate}
  \item Dériver $f(x) = 3x^2 - 5x + 1$ et $g(x) = x^3 - 12x$
        (\cref{thm:g11:deriv:rules})~; et retrouver
        $\left(x^2\right)'$ en $x = 1$ à partir de la définition
        (taux de variation, puis $h \to 0$).
  \item Déterminer la \omterm{def:g11:deriv:tangent}{tangente} à $y = x^3 - 12x$ en $a = 2$.
        Qu'a-t-elle de particulier, et qu'est-ce que cela signale
        (\cref{prop:g11:deriv:extrema})~?
  \item Établir l'\omterm{def:g10:algebra:equation}{équation} générale de la \omterm{def:g11:deriv:tangent}{tangente} à la \omterm{def:g11:quad:parabola}{parabole}
        $y = x^2$ au point $(a, a^2)$~:
        \[
        y = 2ax - a^2 .
        \]
  \item Où cette \omterm{def:g11:deriv:tangent}{tangente} coupe-t-elle l'axe des ordonnées~? En
        déduire la recette du dessinateur~: pour tracer la \omterm{def:g11:deriv:tangent}{tangente}
        en un point $P$ de hauteur $h$ sur la \omterm{def:g11:quad:parabola}{parabole}, joindre $P$
        au point de l'axe situé à la profondeur\,\dots\ sous
        l'origine. L'énoncer.
  \item Construire le \omterm{def:g10:functions:table}{tableau de variations} complet de
        $g(x) = x^3 - 12x$ (signe de $g'$, extremums locaux et leurs
        valeurs, \cref{met:g11:deriv:study}).
\end{enumerate}

\medskip
\noindent\textbf{Partie II --- Le théorème du miroir.}
Soit $P(a, a^2)$ un point de la \omterm{def:g11:quad:parabola}{parabole} $y = x^2$ (avec
$a \neq 0$), $F\left(0, \frac14\right)$ le \omterm{pb:g11:quad:1}{foyer}, et $T$ le point où
la \omterm{def:g11:deriv:tangent}{tangente} en $P$ coupe l'axe des ordonnées (question 4). On
rappelle du \cref{pb:g11:quad:1} que $PF = a^2 + \frac14$.
\begin{enumerate}[resume]
  \item Calculer la distance $FT$.
  \item En déduire que le triangle $FPT$ est isocèle en $F$. Quels
        sont donc les deux angles égaux~?
  \item Le rayon incident qui passe par $P$, parallèle à l'axe, est
        parallèle à la droite $(TF)$ (tous deux verticaux). À l'aide
        des angles alternes-internes, montrer que la \omterm{def:g11:deriv:tangent}{tangente} en $P$
        fait des \emph{angles égaux} avec le rayon vertical incident
        et avec le segment $[PF]$.
  \item Physique~: la lumière se réfléchit sur une courbe comme sur
        sa \omterm{def:g11:deriv:tangent}{tangente}, en repartant sous le même angle qu'à l'arrivée
        (loi de la réflexion). En déduire le \emph{théorème du
        miroir}~: tout rayon parallèle à l'axe se réfléchit en
        passant par le \omterm{pb:g11:quad:1}{foyer} --- et, en renversant le temps, une
        ampoule placée au \omterm{pb:g11:quad:1}{foyer} envoie un faisceau parallèle.
        L'antenne et le phare du \cref{pb:g11:quad:1} sont
        désormais des théorèmes.
  \item Vérifier numériquement le triangle isocèle pour $a = 1$~:
        calculer $T$, $FP$ et $FT$.
\end{enumerate}

\medskip
\noindent\textbf{Partie III --- La machine de Newton.}
Pour résoudre $f(x) = 0$, Newton proposait ceci~: partir d'une
estimation $x_0$, suivre la \omterm{def:g11:deriv:tangent}{tangente} en $x_0$ jusqu'à l'axe, et
prendre son point d'\omterm{def:g10:numbers:interunion}{intersection} comme estimation suivante.
\begin{enumerate}[resume]
  \item Établir la formule de la méthode~:
        \[
        x_{n + 1} = x_n - \frac{f(x_n)}{f'(x_n)} .
        \]
  \item L'appliquer à $f(x) = x^2 - 2$ à partir de $x_0 = 1$~:
        calculer $x_1$, $x_2$, $x_3$ sous forme de fractions
        exactes. Quelle suite vieille de deux mille ans --- et
        quelle machine du \cref{pb:g10:functions:1} --- vient-on de
        reconstruire~? Démontrer l'identité qui explique la
        coïncidence~:
        \[
        x - \frac{x^2 - 2}{2x} = \frac12\left(x +
        \frac2x\right).
        \]
  \item Résoudre $x^3 = 100$ (le croisement des plans d'épargne du
        \cref{pb:g10:reffunc:1})~: appliquer deux pas de Newton à
        $f(x) = x^3 - 100$ à partir de $x_0 = 5$ (quatre
        décimales), et comparer avec le $\sqrt[3]{100}$ de la
        calculatrice.
  \item Sabotage~: essayer de lancer la méthode sur
        $f(x) = x^3 - 12x$ en $x_0 = 2$. Qu'est-ce qui échoue, et
        pourquoi (question 2)~? Énoncer la mise en garde pratique.
  \item À la question 12, le nombre de décimales exactes est passé
        d'environ $1$ à $3$ puis à $6$. Comparer avec la dichotomie
        (division d'un \omterm{def:g10:numbers:interval}{intervalle} par deux à chaque pas) et dire en
        une phrase pourquoi calculatrices, récepteurs GPS et
        entraînement des réseaux de neurones fonctionnent tous sur
        l'idée de Newton.
\end{enumerate}

\medskip
\noindent\textbf{Partie IV --- La poutre la plus résistante.}
On scie une poutre rectangulaire de largeur $w$ et de hauteur $d$
dans un tronc rond de diamètre $30$ cm, de sorte que
$w^2 + d^2 = 900$. La rigidité d'une poutre est proportionnelle à
$w d^2$ (la hauteur compte double~: les solives se posent sur
chant~!).
\begin{enumerate}[resume]
  \item Exprimer la rigidité en \omterm{def:g11:func:function}{fonction} de $w$ seul~:
        $S(w) = 900w - w^3$, sur quel \omterm{def:g10:numbers:interval}{intervalle}~?
  \item Dériver, construire le \omterm{def:g10:functions:table}{tableau de variations}, et déterminer
        les valeurs optimales de $w$ et $d$ (valeurs exactes, puis
        au millimètre).
  \item Montrer que les proportions optimales vérifient
        $d = w\sqrt2$ --- le vieux rapport des charpentiers.
        (Anecdote en prime~: partager le diamètre du tronc en trois
        parts égales et élever des perpendiculaires jusqu'au cercle
        construit exactement ce rectangle.)
  \item Le choix naïf est la poutre \emph{carrée} ($w = d$).
        Calculer sa rigidité et celle de la poutre optimale, et
        donner le gain du charpentier en pourcentage.
  \item Pour finir --- les trois métiers de la \omterm{def:g11:deriv:tangent}{tangente}, une phrase
        pour chacun~: géométrie (le théorème du miroir), calcul
        numérique (la machine de Newton), optimisation (la poutre),
        tous mus par le même objet. Et l'annonce~: l'an prochain, la
        convexité --- de quel côté de la \omterm{def:g11:deriv:tangent}{tangente} se tient la courbe
        --- viendra affûter les trois.
\end{enumerate}
\end{problem}
